\section{Planets Cycles}
\label{sec:cses-planets-cycles}

\problemstatement{Each of $n$ planets has one teleporter. For each starting
planet, output how many planets you visit before returning to a planet you
have already seen (i.e.\ the number of teleports until the first repeat).}

\begin{patternbox}
In a functional graph, starting anywhere you walk a tail and then loop a
cycle. The count until the first repeat is $\textit{depth}+\textit{cycle
length}$: the tail steps to reach the cycle plus the cycle's size. So the
same rho decomposition as Planets Queries II answers every node at once.
\end{patternbox}

\subsection*{Tail length plus cycle length}

Peel to identify cycle nodes; for each cycle record its length. A node on a
cycle simply revisits after $\textit{cycleLen}$ steps. A tail node walks
$\textit{depth}$ steps to reach the cycle, then $\textit{cycleLen}$ more to
return to its entry point -- so its answer is
$\textit{depth}+\textit{cycleLen}(\textit{entry})$. Computing $\textit{depth}$
and the cycle each node leads to (in reverse peel order) gives all answers in
one pass.

\begin{ideabox}[Every Walk Is Tail Then Loop]
The first repeat is always the cycle entry, reached again after one full
loop. Hence the visited count is exactly tail length plus cycle length,
independent of where on the cycle you enter.
\end{ideabox}

\subsection*{Algorithm}

\begin{enumerate}
  \item Peel tree nodes; label cycles with their lengths.
  \item Cycle node answer $=\textit{cycleLen}$; tail node answer
        $=\textit{depth}+\textit{cycleLen}(\textit{lead})$.
  \item Output every node's answer.
\end{enumerate}

\subsection*{Dry run}

\dryrun{$\textit{succ}=[2,3,1,3]$: cycle $1\!\to\!2\!\to\!3$ (length $3$),
tail $4\!\to\!3$.}

\begin{itemize}
  \item Nodes $1,2,3$: answer $3$ each. Node $4$: depth $1$ + cycle $3$
        $=4$.
  \item Output $3\ 3\ 3\ 4$.
\end{itemize}

\subsection*{C++ implementation}

\begin{lstlisting}
#include <bits/stdc++.h>
using namespace std;

int main() {
    int n;
    cin >> n;
    vector<int> succ(n + 1), indeg(n + 1, 0);
    for (int i = 1; i <= n; i++) { cin >> succ[i]; indeg[succ[i]]++; }

    vector<int> deg = indeg, order;
    queue<int> pq;
    for (int i = 1; i <= n; i++) if (deg[i] == 0) pq.push(i);
    vector<bool> removed(n + 1, false);
    while (!pq.empty()) {
        int u = pq.front(); pq.pop();
        removed[u] = true; order.push_back(u);
        if (--deg[succ[u]] == 0) pq.push(succ[u]);
    }

    vector<int> depth(n + 1, 0), lead(n + 1, 0), cycLen(n + 1, 0);
    vector<bool> isCyc(n + 1, false);
    int cid = 0;
    for (int i = 1; i <= n; i++) {
        if (removed[i] || isCyc[i]) continue;
        cid++; int len = 0;
        for (int x = i; !isCyc[x]; x = succ[x]) { isCyc[x] = true; lead[x] = cid; len++; }
        cycLen[cid] = len;
    }
    for (int idx = (int)order.size() - 1; idx >= 0; idx--) {
        int u = order[idx];
        depth[u] = depth[succ[u]] + 1;
        lead[u] = lead[succ[u]];
    }

    for (int i = 1; i <= n; i++)
        cout << depth[i] + cycLen[lead[i]] << " ";
    cout << "\n";
    return 0;
}
\end{lstlisting}

\begin{complexitybox}
\textbf{Time Complexity.} $O(n)$ -- peeling plus one reverse pass.
\textbf{Space Complexity.} $O(n)$.
\end{complexitybox}

\begin{takeawaybox}
Steps-until-repeat in a functional graph is tail depth plus cycle length,
computed for every node from one rho decomposition. Recognising the tail-
then-loop structure turns a per-node simulation into a single linear pass.
\end{takeawaybox}
